% GENERATED FILE. Edit canonical lessons, then run npm run generate:books.
% canonical-source-hash: fnv1a64:320593edae09c033
% canonical-lessons: BN-C119-sorbot, BN-C119-halua, BN-C119-biskut, BN-C119-dhuti, BN-C119-panjabi, BN-R119-first-pass-things-in-the-house, BN-R119-second-pass-nature-and-food

\chapter{Sherbet, Halwa, Biscuit, Dhoti, Long shirt}
\label{ch:bn-sherbet-halwa-biscuit-dhoti-long-shirt}
\hlchaptermodality{119}

\begin{chapteropening}
\textbf{By the end of this chapter:} \emph{I can say a sherbet, halwa, a biscuit, a dhoti and a long shirt.}

Complete the last lesson of chapter 119: I can say a sherbet, halwa, a biscuit, a dhoti and a long shirt.
\end{chapteropening}

\section[\texorpdfstring{\bn{শরবত} (sôrbôt)}{শরবত (sôrbôt)}]{\bn{শরবত} (sôrbôt) — a sherbet}

\label{lesson:BN-C119-sorbot}



\begin{quote}
Before the new one: say the Bengali for puffed rice, then the Bengali for coffee.
\end{quote}

\subsection*{You'll want to know: \bn{শরবত}}
\textbf{\bn{শরবত}} — \emph{sôrbôt} — \textquotedblleft{}a sherbet\textquotedblright{}.

\begin{grammarlens}[title={What you've built}]
One more word for this chapter.
\end{grammarlens}

\subsection*{Your turn}
\begin{itemize}
\raggedright
  \item \emph{Say it:} \emph{sôrbôt}
  \item \emph{Say it:} \emph{sôrbôt}, once more
  \item \emph{Say it:} say \emph{sôrbôt}
  \item \emph{Recall:} say the Bengali for puffed rice, then the Bengali for coffee, then say \emph{sôrbôt} again
\end{itemize}

\begin{tcolorbox}[breakable,colback=teal!4,colframe=teal!35!black,title={Before you move on}]
What does \bn{শরবত} mean? (\textquotedblleft{}A sherbet\textquotedblright{}.) Say it once more.
\end{tcolorbox}
\section[\texorpdfstring{\bn{হালুয়া} (hālua)}{হালুয়া (hālua)}]{\bn{হালুয়া} (hālua) — halwa}

\label{lesson:BN-C119-halua}



\begin{quote}
Before the new one: say the Bengali for coffee, then the Bengali for a sherbet.
\end{quote}

\subsection*{You'll want to know: \bn{হালুয়া}}
\textbf{\bn{হালুয়া}} — \emph{hālua} — \textquotedblleft{}halwa\textquotedblright{}.

\begin{grammarlens}[title={What you've built}]
One more word for this chapter.
\end{grammarlens}

\subsection*{Your turn}
\begin{itemize}
\raggedright
  \item \emph{Say it:} \emph{hālua}
  \item \emph{Say it:} \emph{hālua}, once more
  \item \emph{Say it:} say \emph{hālua}
  \item \emph{Recall:} say the Bengali for coffee, then the Bengali for a sherbet, then say \emph{hālua} again
\end{itemize}

\begin{tcolorbox}[breakable,colback=teal!4,colframe=teal!35!black,title={Before you move on}]
What does \bn{হালুয়া} mean? (\textquotedblleft{}Halwa\textquotedblright{}.) Say it once more.
\end{tcolorbox}
\section[\texorpdfstring{\bn{বিস্কুট} (biskut)}{বিস্কুট (biskut)}]{\bn{বিস্কুট} (biskut) — a biscuit}

\label{lesson:BN-C119-biskut}



\begin{quote}
Before the new one: say the Bengali for a sherbet, then the Bengali for halwa.
\end{quote}

\subsection*{You'll want to know: \bn{বিস্কুট}}
\textbf{\bn{বিস্কুট}} — \emph{biskut} — \textquotedblleft{}a biscuit\textquotedblright{}.

\begin{grammarlens}[title={What you've built}]
One more word for this chapter.
\end{grammarlens}

\subsection*{Your turn}
\begin{itemize}
\raggedright
  \item \emph{Say it:} \emph{biskut}
  \item \emph{Say it:} \emph{biskut}, once more
  \item \emph{Say it:} say \emph{biskut}
  \item \emph{Recall:} say the Bengali for a sherbet, then the Bengali for halwa, then say \emph{biskut} again
\end{itemize}

\begin{tcolorbox}[breakable,colback=teal!4,colframe=teal!35!black,title={Before you move on}]
What does \bn{বিস্কুট} mean? (\textquotedblleft{}A biscuit\textquotedblright{}.) Say it once more.
\end{tcolorbox}
\section[\texorpdfstring{\bn{ধুতি} (dhuti)}{ধুতি (dhuti)}]{\bn{ধুতি} (dhuti) — a dhoti}

\label{lesson:BN-C119-dhuti}



\begin{quote}
Before the new one: say the Bengali for halwa, then the Bengali for a biscuit.
\end{quote}

\subsection*{You'll want to know: \bn{ধুতি}}
\textbf{\bn{ধুতি}} — \emph{dhuti} — \textquotedblleft{}a dhoti\textquotedblright{}.

\begin{grammarlens}[title={What you've built}]
One more word for this chapter.
\end{grammarlens}

\subsection*{Your turn}
\begin{itemize}
\raggedright
  \item \emph{Say it:} \emph{dhuti}
  \item \emph{Say it:} \emph{dhuti}, once more
  \item \emph{Say it:} say \emph{dhuti}
  \item \emph{Recall:} say the Bengali for halwa, then the Bengali for a biscuit, then say \emph{dhuti} again
\end{itemize}

\begin{tcolorbox}[breakable,colback=teal!4,colframe=teal!35!black,title={Before you move on}]
What does \bn{ধুতি} mean? (\textquotedblleft{}A dhoti\textquotedblright{}.) Say it once more.
\end{tcolorbox}
\section[\texorpdfstring{\bn{পাঞ্জাবি} (pānjābi)}{পাঞ্জাবি (pānjābi)}]{\bn{পাঞ্জাবি} (pānjābi) — a long shirt}

\label{lesson:BN-C119-panjabi}



\begin{quote}
Before the new one: say the Bengali for a biscuit, then the Bengali for a dhoti.
\end{quote}

\subsection*{You'll want to know: \bn{পাঞ্জাবি}}
\textbf{\bn{পাঞ্জাবি}} — \emph{pānjābi} — \textquotedblleft{}a long shirt\textquotedblright{}.

\begin{grammarlens}[title={What you've built}]
One more word for this chapter.
\end{grammarlens}

\subsection*{Your turn}
\begin{itemize}
\raggedright
  \item \emph{Say it:} \emph{pānjābi}
  \item \emph{Say it:} \emph{pānjābi}, once more
  \item \emph{Say it:} say \emph{pānjābi}
  \item \emph{Recall:} say the Bengali for a biscuit, then the Bengali for a dhoti, then say \emph{pānjābi} again
\end{itemize}

\begin{tcolorbox}[breakable,colback=teal!4,colframe=teal!35!black,title={Before you move on}]
What does \bn{পাঞ্জাবি} mean? (\textquotedblleft{}A long shirt\textquotedblright{}.) Say it once more.
\end{tcolorbox}
\section[(dialogue)]{Things in the house — a first pass}

\label{lesson:BN-R119-first-pass-things-in-the-house}



\begin{quote}
Say the words for a dhoti and a long shirt.
\end{quote}

\subsection*{Your turn}
Say these aloud:

\begin{itemize}
\raggedright
  \item scissors, a hurricane lamp, a candle, matches, coal
  \item a rope, a nail, a hammer, an axe, an umbrella
  \item a box, a shelf, a wall, a veranda, a garden
  \item a well, mud, dust, thunder, a rainbow
  \item fog, dew, a root, a seed, a thorn
\end{itemize}

\begin{tcolorbox}[breakable,colback=teal!4,colframe=teal!35!black,title={Before you move on}]
Scissors? (\textbf{\bn{কাঁচি}}, \emph{kā̃chi}.) A hurricane lamp? (\textbf{\bn{হারিকেন}}, \emph{hāriken}.) A candle? (\textbf{\bn{মোমবাতি}}, \emph{mombāti}.) Matches? (\textbf{\bn{দেশলাই}}, \emph{deshlāi}.) Coal? (\textbf{\bn{কয়লা}}, \emph{kôylā}.) A rope? (\textbf{\bn{দড়ি}}, \emph{dôṛi}.) A nail? (\textbf{\bn{পেরেক}}, \emph{perek}.) A hammer? (\textbf{\bn{হাতুড়ি}}, \emph{hātuṛi}.) An axe? (\textbf{\bn{কুড়ুল}}, \emph{kuṛul}.) An umbrella? (\textbf{\bn{ছাতা}}, \emph{chhātā}.) A box? (\textbf{\bn{বাক্স}}, \emph{bāksô}.) A shelf? (\textbf{\bn{তাক}}, \emph{tāk}.) A wall? (\textbf{\bn{দেওয়াল}}, \emph{deoyāl}.) A veranda? (\textbf{\bn{বারান্দা}}, \emph{bārāndā}.) A garden? (\textbf{\bn{বাগান}}, \emph{bāgān}.) A well? (\textbf{\bn{কুয়ো}}, \emph{kuyo}.) Mud? (\textbf{\bn{কাদা}}, \emph{kādā}.) Dust? (\textbf{\bn{ধুলো}}, \emph{dhulo}.) Thunder? (\textbf{\bn{বজ্র}}, \emph{bôjrô}.) A rainbow? (\textbf{\bn{রংধনু}}, \emph{rôngdhônu}.) Fog? (\textbf{\bn{কুয়াশা}}, \emph{kuyāshā}.) Dew? (\textbf{\bn{শিশির}}, \emph{shishir}.) A root? (\textbf{\bn{শিকড়}}, \emph{shikôṛ}.) A seed? (\textbf{\bn{বীজ}}, \emph{bij}.) A thorn? (\textbf{\bn{কাঁটা}}, \emph{kā̃ṭā}.)
\end{tcolorbox}
\section[(dialogue)]{Nature and food — a second pass}

\label{lesson:BN-R119-second-pass-nature-and-food}



\begin{quote}
Say the words for bamboo and a mango tree.
\end{quote}

\subsection*{Your turn}
Say these aloud:

\begin{itemize}
\raggedright
  \item bamboo, a mango tree, a coconut, a pineapple, a papaya
  \item a watermelon, an aubergine, garlic, ginger, a tomato
  \item beans, a pumpkin, leafy greens, a radish, uncooked rice
  \item rasgulla, rice pudding, a luchi, puffed rice, coffee
  \item a sherbet, halwa, a biscuit, a dhoti, a long shirt
\end{itemize}

\begin{tcolorbox}[breakable,colback=teal!4,colframe=teal!35!black,title={Before you move on}]
Bamboo? (\textbf{\bn{বাঁশ}}, \emph{bā̃sh}.) A mango tree? (\textbf{\bn{আমগাছ}}, \emph{āmgāchh}.) A coconut? (\textbf{\bn{নারকেল}}, \emph{nārkel}.) A pineapple? (\textbf{\bn{আনারস}}, \emph{ānārôs}.) A papaya? (\textbf{\bn{পেঁপে}}, \emph{pepe}.) A watermelon? (\textbf{\bn{তরমুজ}}, \emph{tôrmuj}.) An aubergine? (\textbf{\bn{বেগুন}}, \emph{begun}.) Garlic? (\textbf{\bn{রসুন}}, \emph{rôsun}.) Ginger? (\textbf{\bn{আদা}}, \emph{ādā}.) A tomato? (\textbf{\bn{টমেটো}}, \emph{tômeṭo}.) Beans? (\textbf{\bn{শিম}}, \emph{shim}.) A pumpkin? (\textbf{\bn{কুমড়ো}}, \emph{kumṛo}.) Leafy greens? (\textbf{\bn{শাক}}, \emph{shāk}.) A radish? (\textbf{\bn{মুলো}}, \emph{mulo}.) Uncooked rice? (\textbf{\bn{চাল}}, \emph{chāl}.) Rasgulla? (\textbf{\bn{রসগোল্লা}}, \emph{rôsôgollā}.) Rice pudding? (\textbf{\bn{পায়েস}}, \emph{pāyes}.) A luchi? (\textbf{\bn{লুচি}}, \emph{luchi}.) Puffed rice? (\textbf{\bn{মুড়ি}}, \emph{muṛi}.) Coffee? (\textbf{\bn{কফি}}, \emph{kôphi}.) A sherbet? (\textbf{\bn{শরবত}}, \emph{sôrbôt}.) Halwa? (\textbf{\bn{হালুয়া}}, \emph{hālua}.) A biscuit? (\textbf{\bn{বিস্কুট}}, \emph{biskut}.) A dhoti? (\textbf{\bn{ধুতি}}, \emph{dhuti}.) A long shirt? (\textbf{\bn{পাঞ্জাবি}}, \emph{pānjābi}.)
\end{tcolorbox}
